\section{Character Spectra and Moment Profiles}

This appendix collects the exact character-operator spectra and the
walk-moment data used in the main text.

\subsection{Character conventions}

The native Klein group is
\[
V_4^{\mathrm{nat}}
=
\{1,a,b,ab\}.
\]

Its real character group is
\[
\widehat{V_4}
=
\{\chi_1,\chi_a,\chi_b,\chi_{ab}\}.
\]

The nontrivial characters are indexed by their kernels:
\[
\ker\chi_a=\langle a\rangle,
\]
\[
\ker\chi_b=\langle b\rangle,
\]
and
\[
\ker\chi_{ab}=\langle ab\rangle.
\]

Each character sector has dimension fifteen.

\subsection{Signed quotient operators}

For a selected chamber representative system and quotient voltage
assignment
\[
\omega:E(L(P))\longrightarrow V_4,
\]
the signed quotient adjacency operator is
\[
(A_\chi)_{ij}
=
\begin{cases}
\chi(\omega_{ij}), & ij\in E(L(P)),\\
0, & \text{otherwise}.
\end{cases}
\]

The corresponding character Laplacian is
\[
L_\chi=4I-A_\chi.
\]

For the trivial character,
\[
A_{\chi_1}=A_{L(P)}.
\]

\subsection{Gauge covariance}

Under the chamber-representative change
\[
s_i'=h_is_i,
\qquad
h_i\in V_4,
\]
define
\[
D_\chi
=
\operatorname{diag}
\bigl(
\chi(h_0),\ldots,\chi(h_{14})
\bigr).
\]

The switched operators satisfy
\[
A_\chi'
=
D_\chi A_\chi D_\chi
\]
and
\[
L_\chi'
=
D_\chi L_\chi D_\chi.
\]

Thus every character spectrum is independent of the selected quotient
gauge.

The tested controls produce fifty-five distinct raw voltage profiles,
but no character-spectrum changes.

\subsection{Trivial-character spectrum}

The trivial-sector Laplacian spectrum is
\[
\operatorname{Spec}(L_{\chi_1})
=
\{0^1,2^5,5^4,6^5\}.
\]

Equivalently, the quotient adjacency spectrum is
\[
\operatorname{Spec}(A_{\chi_1})
=
\{4^1,2^5,(-1)^4,(-2)^5\}.
\]

This is the spectrum of
\[
L(P).
\]

\subsection{The \(a\)-sector}

The distinguished nontrivial sector has Laplacian spectrum
\[
\operatorname{Spec}(L_{\chi_a})
=
\{1^4,4^5,6^6\}.
\]

The corresponding adjacency spectrum is
\[
\operatorname{Spec}(A_{\chi_a})
=
\{3^4,0^5,(-2)^6\}.
\]

This sector is fixed under the full automorphism group because every
automorphism fixes the native involution \(a\).

\subsection{The \(b\)- and \(ab\)-sectors}

The exchanged sectors have equal Laplacian spectra:
\[
\operatorname{Spec}(L_{\chi_b})
=
\operatorname{Spec}(L_{\chi_{ab}}).
\]

Their common spectrum is
\[
\left\{
\left(\frac{3-\sqrt5}{2}\right)^3,
3^4,
\left(\frac{3+\sqrt5}{2}\right)^3,
6^5
\right\}.
\]

The corresponding adjacency spectrum is
\[
\left\{
\left(\frac{5+\sqrt5}{2}\right)^3,
1^4,
\left(\frac{5-\sqrt5}{2}\right)^3,
(-2)^5
\right\}.
\]

The selected signed matrices
\[
A_{\chi_b}
\quad\text{and}\quad
A_{\chi_{ab}}
\]
need not be equal. Their isospectrality is enforced by native
automorphisms exchanging
\[
b\longleftrightarrow ab.
\]

\subsection{Full sector summary}

The Laplacian spectra are summarized by

\[
\begin{array}{c|l}
\text{sector} & \text{Laplacian spectrum}\\
\hline
\chi_1 &
\{0^1,2^5,5^4,6^5\}\\
\chi_a &
\{1^4,4^5,6^6\}\\
\chi_b &
\left\{
\left(\frac{3-\sqrt5}{2}\right)^3,
3^4,
\left(\frac{3+\sqrt5}{2}\right)^3,
6^5
\right\}\\
\chi_{ab} &
\left\{
\left(\frac{3-\sqrt5}{2}\right)^3,
3^4,
\left(\frac{3+\sqrt5}{2}\right)^3,
6^5
\right\}
\end{array}
\]

Each row has total multiplicity fifteen.

The four rows together have total multiplicity
\[
4\cdot15=60.
\]

\subsection{Registered adjacency moments}

For every registered ten-cycle \(C\),
\[
m_k^A(C)
=
x_C^{\mathsf T}A^kx_C.
\]

The exact profile through order eight is

\[
\begin{array}{c|rrrrrrrrr}
k & 0 & 1 & 2 & 3 & 4 & 5 & 6 & 7 & 8\\
\hline
m_k^A
&
10
&
20
&
80
&
240
&
880
&
3040
&
11200
&
40960
&
153600
\end{array}
\]

The fourth-order value
\[
m_4^A=880
\]
selects exactly the twenty-four registered cycles among all \(1824\)
simple ten-cycles.

\subsection{Registered Laplacian moments}

For every registered ten-cycle,
\[
m_k^L(C)
=
x_C^{\mathsf T}L^kx_C.
\]

The exact profile through order eight is

\[
\begin{array}{c|rrrrrrrrr}
k & 0 & 1 & 2 & 3 & 4 & 5 & 6 & 7 & 8\\
\hline
m_k^L
&
10
&
20
&
80
&
400
&
2160
&
12000
&
67520
&
382720
&
2181120
\end{array}
\]

The fourth-order value
\[
m_4^L=2160
\]
is equivalent to the adjacency selector because
\[
L=4I-A.
\]

\subsection{Registered and frontier comparison}

The registered and frontier cycles agree in the tested adjacency
moments through order three.

At order four,
\[
m_4^A(\text{registered})=880,
\]
while
\[
m_4^A(\text{frontier})=896.
\]

The difference is
\[
16.
\]

The endpoint-distance comparison is

\[
\begin{array}{c|rr}
 & \text{registered} & \text{frontier}\\
\hline
d_4 & 20 & 28\\
d_5 & 0 & 8
\end{array}
\]

Thus the frontier surplus consists exactly of
\[
8
\]
distance-four external returns and
\[
8
\]
distance-five antipodal returns.

\subsection{Boundary-crossing profile}

For a registered cycle, the four-step boundary-crossing counts are

\[
\begin{array}{c|rrr}
\text{crossings} & 0 & 2 & 4\\
\hline
\text{count} & 160 & 560 & 160
\end{array}
\]

The total is
\[
160+560+160=880.
\]

The value
\[
560
\]
for exactly two boundary crossings is itself an exact registered-cycle
selector in the full ten-cycle census.

\subsection{Exact scalar selectors}

The fourth-order anatomy audit records several exact scalar selectors
for the registered class:

\[
\text{boundary-crossing count at }2=560,
\]
\[
\text{endpoint distance }4=20,
\]
\[
\text{endpoint distance }5=0,
\]
\[
\text{three-outside-vertex count}=80,
\]
and
\[
\text{five-distinct-quotient count}=100.
\]

These selectors are equivalent only within the complete audited
ten-cycle census. No claim is made that the same scalar values select
an analogous class in arbitrary graphs.

\subsection{Machine-readable sources}

The authoritative spectral sources are:

\begin{enumerate}
    \item
    \texttt{artifacts/json/native\_g60\_v4\_character\_energy\_audit\_053.json};

    \item
    \texttt{artifacts/json/native\_g60\_character\_laplacian\_closure\_audit\_054.json};

    \item
    \texttt{artifacts/json/native\_g60\_character\_exchange\_symmetry\_audit\_055.json};

    \item
    \texttt{artifacts/json/native\_g60\_character\_exchange\_c2\_theorem\_056.json};

    \item
    \texttt{artifacts/json/native\_g60\_v4\_voltage\_character\_quotients\_057.json};

    \item
    \texttt{artifacts/json/native\_g60\_v4\_gauge\_switching\_audit\_058.json};

    \item
    \texttt{artifacts/json/native\_g60\_registered\_cycle\_dynamical\_signature\_audit\_064.json};

    \item
    \texttt{artifacts/json/native\_g60\_registered\_fourth\_moment\_anatomy\_audit\_065.json};

    \item
    \texttt{artifacts/json/native\_g60\_registered\_antipodal\_exclusion\_audit\_066.json}.
\end{enumerate}

The irrational spectral values are exact algebraic numbers arising
from finite matrices. No physical meaning is assigned to them in this
paper.
